% Plano do TCC em linguagem simples. Fonte editavel do Plano_TCC_simples.pdf.
% Compilar:  pdflatex Plano_TCC_simples.tex   (rodar duas vezes se mudar o numero de paginas)
\documentclass[10pt,a4paper]{article}

\usepackage[utf8]{inputenc}
\usepackage[T1]{fontenc}
\usepackage[brazilian]{babel}
\usepackage[scaled=0.95]{helvet}
\renewcommand{\familydefault}{\sfdefault}
\usepackage[a4paper,left=2cm,right=2cm,top=1.5cm,bottom=1.8cm,footskip=0.8cm]{geometry}
\usepackage[table]{xcolor}
\usepackage{array,tabularx,enumitem,fancyhdr,titlesec,needspace}
\usepackage[hidelinks,pdftitle={Plano do TCC: biblioteca prossumidor},
            pdfauthor={João Carlos Jachenski Junior}]{hyperref}

% ------------------------------------------------------------------ cores
\definecolor{acento}{HTML}{0F5C6E}
\definecolor{claro}{HTML}{E8F1F3}
\definecolor{cinza}{HTML}{5A6470}
\definecolor{linha}{HTML}{C9D3D8}
\definecolor{zebra}{HTML}{F5F8F9}
\definecolor{escrita}{HTML}{9FC3CC}
\definecolor{opcional}{HTML}{C9D3D8}
\definecolor{texto}{HTML}{1D232A}
\color{texto}

% ------------------------------------------------------------------ aparencia
\setlength{\parindent}{0pt}
\setlength{\parskip}{4pt}
\linespread{1.08}
\titleformat{\section}{\color{acento}\bfseries\large}{\thesection.}{0.4em}{}
\titlespacing*{\section}{0pt}{10pt}{4pt}
\setlist[itemize]{leftmargin=13pt,itemsep=2pt,topsep=2pt,label=\textbullet}
\pagestyle{fancy}
\fancyhf{}
\renewcommand{\headrulewidth}{0pt}
\fancyfoot[L]{\scriptsize\color{cinza}prossumidor -- plano do TCC em linguagem simples}
\fancyfoot[R]{\scriptsize\color{cinza}página \thepage}

% Cabecalho de tabela: fundo da cor de destaque, texto branco e negrito.
\newcommand{\cab}[1]{\cellcolor{acento}\color{white}\bfseries #1}
\newcolumntype{Y}{>{\raggedright\arraybackslash}X}
\renewcommand{\arraystretch}{1.35}
\newcommand{\nota}[1]{{\small\itshape\color{cinza}#1\par}}
% Caixa de destaque com filete a esquerda.
\newcommand{\caixa}[1]{%
  \noindent\begin{tabular}{!{\color{acento}\vrule width 3pt}>{\columncolor{claro}}p{\dimexpr\linewidth-3pt-2\tabcolsep}}
  \rule{0pt}{14pt}#1\vspace{6pt}\\
  \end{tabular}\par\vspace{4pt}}
% Celula do cronograma: mes ocupado.
\newcommand{\on}{\cellcolor{acento}}
\newcommand{\es}{\cellcolor{escrita}}
\newcommand{\op}{\cellcolor{opcional}}

\begin{document}
\small

% ================================================================== titulo
{\color{acento}\bfseries\LARGE Plano do TCC: biblioteca prossumidor\par}
\vspace{2pt}
{\color{cinza}O que é, o que já foi feito e o que falta \;|\; João Carlos Jachenski Junior \;|\; Lactec \;|\; revisão de 09/10/2026\par}
\vspace{8pt}

\caixa{\textbf{Em uma frase:} o TCC é uma biblioteca em Python que prevê quanta energia uma casa ou
usina solar vai \textbf{gerar} e \textbf{consumir} no dia seguinte, desde a leitura do medidor até um
modelo pequeno que roda num aparelho barato (Raspberry Pi ou ESP32).}

% ================================================================== proposta
\section{A proposta}
Hoje o projeto tem dezenas de scripts soltos que repetem as mesmas contas (datas de teste, métricas,
limpeza de dados). Quando uma coisa muda, é preciso corrigir à mão em muitos lugares, e já houve erros
por isso. A ideia é juntar tudo numa \textbf{biblioteca única, testada e reaproveitável}. Ela não
reinventa o que já existe: usa bibliotecas prontas (pvlib, LightGBM, modelos de IA públicos) e só
acrescenta o que é próprio deste projeto.

\rowcolors{2}{white}{zebra}
\begin{tabularx}{\linewidth}{@{}>{\raggedright\arraybackslash}p{3.2cm}Y@{}}
\cab{Parte da biblioteca} & \cab{O que faz, em poucas palavras} \\
io          & Lê os arquivos reais: inversor Fronius, medidores das casas, dados de clima. \\
qualidade   & Acha dados ruins: buracos, noite, troca de horário, casa sem painel solar. \\
fisica      & Calcula a posição do sol e quanta luz chegaria num dia sem nuvens. \\
protocolo   & Regras únicas de teste: que dias prever, como comparar, como medir o erro. \\
previsao    & Um jeito único de usar qualquer modelo, do mais simples aos modelos de IA. \\
diagnostico & Descobre painel sujo ou com defeito usando só os dados do inversor. \\
embarcado   & Encolhe o modelo para caber e rodar num aparelho pequeno. \\
economia    & Calcula a conta de luz: tarifas da ANEEL, bandeiras e a Lei 14.300. \\
\end{tabularx}

% ================================================================== feito
\section{O que já foi feito}
\begin{itemize}
  \item \textbf{Teste grande de modelos de previsão.} \textbf{51 modelos} em 6 testes para o artigo (usinas,
        geração FV das casas e consumo das casas, cada um sem e com o clima do dia) e mais o regime com a
        previsão do tempo real nas usinas: 243 resultados. Cada modelo é retreinado a cada período de teste,
        e as diferenças são checadas com testes estatísticos. Entraram também os clássicos (ARIMA, ETS,
        Theta, MSTL, Ridge, Random Forest, MLP, k-NN).
  \item \textbf{As 20 casas foram identificadas} (Dingle, Irlanda, domínio público) e ganharam o clima
        medido da estação local. O relógio delas (hora local com horário de verão) foi descoberto
        comparando a geração solar com a posição do sol.
  \item \textbf{Autorização} da gestão para uso acadêmico (05/10/2026) -- etapa E0.
  \item \textbf{Base da biblioteca} (etapa E1): pacote instalável, checagem dos dados de entrada (com
        horário de verão), dados de exemplo, \textbf{304 testes automáticos} no GitHub, incluindo um que
        recalcula os 243 resultados publicados e exige que batam exatamente.
  \item \textbf{Revisão de 09/10/2026:} erros de método achados e corrigidos (código da casa tratado como
        número em três modelos; planilhas geradas no meio de uma rodada), e \textbf{testes-piloto} que
        decidiram dois comportamentos da biblioteca (seção 4).
\end{itemize}

% ================================================================== dados
\section{Os dados}
\rowcolors{2}{white}{zebra}
\begin{tabularx}{\linewidth}{@{}>{\raggedright\arraybackslash}p{2.4cm}Y>{\raggedright\arraybackslash}p{4.6cm}@{}}
\cab{Dados} & \cab{Origem} & \cab{Pode publicar?} \\
2 usinas solares & Cascavel (3,96 kWp) e Curitiba ($\sim$9,6 kWp), de proprietários particulares.
                 & \textbf{Não.} Sem consentimento. Nomes trocados pela cidade. \\
20 casas & Comunidade de energia em Dingle, Irlanda, 2020, medição a cada minuto, com estação
           meteorológica. Khadem et al. (2024), \textit{Scientific Data}, doi:10.1038/s41597-024-03454-2.
         & \textbf{Sim.} Domínio público (CC0). Basta citar o artigo. \\
\end{tabularx}

% ================================================================== licoes
\needspace{8cm}
\section{O que os testes ensinaram para a biblioteca}
\begin{itemize}
  \item \textbf{O clima pesa cerca de 10 vezes mais que o modelo.} Nas usinas, dar o clima ao modelo
        sobe o skill em +0,46 (mediana); trocar o 1º pelo 5º modelo muda 0,04. Por isso o maior
        investimento vai para \textbf{obter clima e previsão do tempo melhores}, não para mais modelos.
  \item \textbf{A mediana de todos os modelos empata com o melhor nos 6 testes.} Vira a saída padrão.
        \textit{Testado:} escolher o melhor da rodada anterior piora em todos os casos (até $-$0,23),
        e ficar só com os 5 ou 10 melhores não é estável. A biblioteca \textbf{não} escolhe modelo.
  \item \textbf{Correção de viés só no consumo.} \textit{Testado:} os modelos de IA preveem $\sim$20\% a
        menos de consumo; corrigir pelo viés dos últimos 14 dias leva o erro médio a $<$1\% e sobe a
        mediana de +0,191 para +0,207 (empata com o melhor). Na geração solar a mesma correção
        \textbf{piora} até $-$0,10: lá ela fica desligada.
  \item \textbf{O consumo de uma casa, hora a hora, é em grande parte imprevisível} (R² 0,37); somando as
        20 casas sobe para 0,64, e o total do dia para 0,58. A biblioteca mede o erro em várias
        resoluções (hora, blocos de tarifa, pico da noite, dia), para cada uso escolher a sua. No total do
        dia, a média dos últimos 7 dias quase empata com os modelos.
  \item \textbf{Medir, não supor.} O fuso das casas estava errado e só apareceu comparando a geração com o
        sol. Isso vira uma verificação automática do relógio no módulo de qualidade.
  \item \textbf{Cada modelo declara o que exige}: GPU, licença dos pesos e se envia dados para fora (um
        deles roda em servidor pago; dados privados nunca vão para lá sem permissão).
\end{itemize}

% ================================================================== idioma
\needspace{6cm}
\section{Decisão pendente: biblioteca nacional ou internacional}
\caixa{\textbf{Decidir até o fim de outubro de 2026, antes do grosso da etapa E2.} Trocar depois custa
várias vezes mais, porque cada módulo escrito teria de ser traduzido.}

\rowcolors{2}{white}{zebra}
\begin{tabularx}{\linewidth}{@{}>{\raggedright\arraybackslash}p{3.0cm}YY@{}}
\cab{} & \cab{Nacional (\texttt{prossumidor}, em português)} & \cab{Internacional (\texttt{prosumer}, em inglês)} \\
Para quem & distribuidoras, integradores e pesquisadores no Brasil & qualquer pesquisador ou empresa \\
Destaque & tarifas da ANEEL, Lei 14.300, Copel: ninguém mais faz & a parte de previsão e o embarcado \\
Custo extra & nenhum & $\sim$10 h agora (nomes e código) + $\sim$20 h de documentação em inglês \\
Publicação & SoftwareX ou revista nacional & JOSS ou SoftwareX \\
\end{tabularx}

Três perguntas decidem: (1) quem vai usar primeiro, a Lactec e o setor elétrico brasileiro ou a
comunidade científica? (2) há plano de mestrado fora do Brasil? (3) a licença que a Lactec decidir
permite publicar aberto? Os dois caminhos usam o mesmo plano abaixo; só muda a E6.

% ================================================================== plano
\needspace{8cm}
\section{O plano: testar antes de investir}
Toda melhoria começa por um \textbf{piloto curto} (1 a 4 horas, num caso só). Só vira trabalho completo
se passar no critério da coluna ``Portão''. Assim nenhuma etapa longa é feita às cegas. Horas com
cerca de 15 horas por semana.

\rowcolors{2}{white}{zebra}
\begin{tabularx}{\linewidth}{@{}>{\raggedright\arraybackslash}p{2.7cm}Y>{\raggedright\arraybackslash}p{4.0cm}>{\raggedleft\arraybackslash}p{0.9cm}@{}}
\cab{Etapa} & \cab{O que é} & \cab{Portão (piloto)} & \cab{Horas} \\
F0 -- Decisões & Idioma da biblioteca; licença (Lactec); comprar Raspberry Pi 5 e ESP32. & -- & 0 \\
F1 -- Fechar o teste & Congelar os ambientes; avaliação por resolução (hora, tarifa, pico, dia);
  correção de viés do consumo (já testada); ARIMA, MLP e 3 modelos novos que valem a pena
  (TimesFM-3, FlowState, BuildingsBench). & ARIMA: 1 caso em menos de 2 h; MLP: ganho $\geq$0,05;
  modelos novos: entrar no grupo da frente num caso. & 17 \\
F2 -- Pilotos de melhoria & (a) combinar 2--3 previsões do tempo; (b) estimar a inclinação e a direção
  de cada telhado pelos dados; (c) prever uma faixa (quantis); (d) atualizar a previsão ao meio-dia;
  (e) ajuste fino de um modelo de IA nos nossos dados. & cada um: skill $+$0,02 ou mais num caso, ou
  faixa bem calibrada (c). & 20 \\
F2b -- Levar para a biblioteca & Só os pilotos que passaram viram funções da biblioteca. & estimativa & 40 \\
E2 -- Versão 0.1 & Núcleo: ler e limpar dados (com o relógio verificado pelo sol), física do sol,
  regras de teste em várias resoluções, ranking e empate, modelos com o que exigem declarado, mediana
  como saída, dados reais de exemplo (Dingle). & testes automáticos; cobertura $\geq$80\% & 120 \\
E3 -- Checagem 1 & Refazer o teste grande usando só a biblioteca. & os 243 resultados batem & 30 \\
E4 -- Versão 0.2 & Diagnóstico de painel sujo ou com defeito; modelos pequenos para aparelhos
  (TTM, SparseTSF, Ridge, Random Forest; previsão do dia e por bloco de tarifa). & diagnóstico de
  Cascavel reproduzido & 130 \\
E5 -- Checagem 2 & Medir tamanho, acerto, velocidade e energia no Raspberry Pi 5 e no ESP32. & tabela
  medida no aparelho & 50 \\
E6 -- Versão 0.3 & Conta de luz (tarifa branca, bandeiras, Lei 14.300), manual e publicação. & licença
  decidida & 80 \\
E7 -- Escrita & O TCC. Começa já, em paralelo. & aprovação do orientador & 80 \\
\rowcolor{claro}\textbf{Total} & & & \textbf{567} \\
\end{tabularx}

\nota{Total de $\sim$38 semanas (9 meses). Se a decisão for internacional: $+$30 h. Subprodutos, depois da
biblioteca e sem atrasar a E2: artigo de congresso com os regimes A e B ($\sim$25 h; a planilha já está
pronta) e artigo de software ($\sim$30 h).}

% ================================================================== cortes
\needspace{6cm}
\section{O que não vale a pena (e por quê)}
\begin{itemize}
  \item \textbf{Escolher automaticamente ``o melhor modelo''}: testado, piora ou é instável.
  \item \textbf{Correção de viés na geração solar}: testado, piora.
  \item \textbf{Os outros 13 modelos da lista de pendentes}: Lag-Llama, MOMENT, VisionTS e Moirai 1 ficam
        atrás dos modelos de IA mais novos já testados nos benchmarks públicos; CycleNet, DUET e
        Crossformer são de famílias já cobertas por 24 redes, e nenhuma rede passou da mediana; Mamba e
        S-Mamba exigem Linux e GPU; LLMTime e Time-LLM são da família LLM, que não ganha dos modelos sem
        LLM e estoura a memória; TimeGPT é pago e envia os dados para fora; TimeGrad só volta se o
        piloto de faixas (F2c) passar.
  \item \textbf{Um modelo separado para o consumo do dia}: a soma das previsões por hora já dá o total
        do dia, e nele os modelos quase não passam da média de 7 dias.
  \item \textbf{Mais modelos em geral}: o clima pesa 10 vezes mais. O tempo vai para clima, física e a
        biblioteca.
\end{itemize}

% ================================================================== cronograma
\needspace{9cm}
\section{Cronograma sugerido}
\rowcolors{1}{}{}
\setlength{\tabcolsep}{0pt}
\renewcommand{\arraystretch}{1.5}
\arrayrulecolor{linha}
\begin{tabular}{|>{\hspace{4pt}}m{3.4cm}|*{11}{>{\centering\arraybackslash}m{1.2cm}|}}
\hline
\cab{Etapa} & \cab{out} & \cab{nov} & \cab{dez} & \cab{jan} & \cab{fev} & \cab{mar} & \cab{abr} & \cab{mai} & \cab{jun} & \cab{jul} & \cab{ago} \\ \hline
F0 Decisões            & \on &     &     &     &     &     &     &     &     &     &     \\ \hline
F1 Fechar o teste      & \on &     &     &     &     &     &     &     &     &     &     \\ \hline
F2 Pilotos             & \on & \on &     &     &     &     &     &     &     &     &     \\ \hline
E2 Versão 0.1          &     & \on & \on & \on &     &     &     &     &     &     &     \\ \hline
E3 Checagem 1          &     &     &     &     & \on &     &     &     &     &     &     \\ \hline
E4 Versão 0.2          &     &     &     &     & \on & \on & \on &     &     &     &     \\ \hline
E5 Hardware real       &     &     &     &     &     &     &     & \on &     &     &     \\ \hline
E6 Versão 0.3          &     &     &     &     &     &     &     &     & \on &     &     \\ \hline
E7 Escrita do TCC      & \es & \es & \es & \es & \es & \es & \es & \es & \es & \es &     \\ \hline
Artigos (subprodutos)  &     & \op & \op &     &     &     &     &     &     &     & \op \\ \hline
\end{tabular}
\setlength{\tabcolsep}{6pt}
\renewcommand{\arraystretch}{1.35}

\vspace{3pt}
\nota{Outubro de 2026 a julho de 2027 para a biblioteca e o TCC. Os artigos (cinza) são subprodutos: o de
congresso aproveita a planilha pronta em nov--dez; o de software, depois da publicação.}

% ================================================================== atencao
\section{O que pode atrasar (resolver logo)}
\begin{itemize}
  \item \textbf{Decisão de idioma} (seção 5) e \textbf{licença} (Lactec): travam a E2 e a E6.
  \item \textbf{Orientador e coorientador} confirmados por escrito, se ainda não estiver feito.
  \item \textbf{Comprar o Raspberry Pi 5 e o ESP32} com antecedência (etapa E5).
  \item \textbf{Memória da máquina de casa}: três rodadas morreram por falta de RAM. Rodar em lotes
        pequenos e não misturar filas de CPU e GPU.
\end{itemize}

\end{document}
